\section{Structural Estimation of DDCs\label{sec:ddc}}

The term ``structural'' refers to the attempt to infer the underlying preferences and beliefs of economic
agents rather than simply summarizing their behavior, the main goal of reduced-form econometrics.
In the MDP framework, an agent's decision rule $\pi(a|s)$ is the ``reduced-form model'' capturing behavior, whereas structural econometrics infers the underlying preferences and beliefs $\{\beta,r,p\}$. Key references include \citet{rust:1987}, \citet{rust:1988}, \citet{rust:1994}, \citet{AguirregabiriaMira:2002},
and the survey by \citet{AguirregabiriaMira:2010}.\footnote{\footnotesize See \citet{rust:2014} for a discussion 
of the origins of SE, dating back to work at the Cowles Foundation in the 1940s and the rationale for this literature including
the argument of \citet{lucas:1976} that structural models are superior to reduced-form for counterfactual prediction of policy changes.
The development of DP in the 1950s and 60s 
revolutionized economics by enabling models of sequential
decision making under uncertainty, dynamic strategic interactions, and dynamic equilibria 
in markets which are key features of nearly all economic problems.  DP lead to new dynamic economic theories
based on {\it rational expectations\/} that
showed how beliefs about uncertain future events affect the current choices of rational forward looking agents. 
This in turn lead to empirical methods for estimating and testing how well these new theories
explain behavior, an area now known as {\it dynamic structural estimation\/} (DSE).} 

An obstacle to inference is the {\it statistical degeneracy\/} of the optimal decision rule $\delta^*(s)$ 
in equation (\ref{eq:bellman_policy}), i.e. the fact that it is generally a pure strategy (i.e. deterministic function of $s$). 
Different individuals in the same observed state $s$ will typically make different choices, so
to enable the  model to explain this heterogeneity we  augment the
state variable as $x=(s,\varepsilon)$  where
$\varepsilon$ is a vector of idiosyncratic preference shocks observed by the agent but not by others.
Even though $\delta^*(x,\varepsilon)$ is still a pure strategy from the agent's standpoint, it produces choices that appear
random from the observer's standpoint. Under additional assumptions that include the assumption $A(s)$ is finite for all $s \in S$,
we can derive a {\it conditional choice probability\/} (CCP) $\pi(a|s)$ implied by the 
strategy $\delta^*$ that admits positive probabilities for all feasible choices $a \in A(s)$. Using this, we can form a likelihood
function to infer $\{\beta,r,p\}$ from data on agents' states and choices.

Building on \citet{mcfadden:1973}'s static discrete choice
models, we can derive a simple expression for the CCP $\pi(a|s)$ known as 
the {\it multinomial logit model\/} or ``softmax'' function in the RL literature.
Suppose that for each $s$, $\varepsilon$ is a vector with the same length as $|A(s)|$ the number of actions, so for
$\varepsilon(a)$ is the component of $\varepsilon$ associated with action $a \in A(s)$. Assume that 
the reward for action $a$ in state $(s,\varepsilon)$ can be written as $r(s,a)+\varepsilon(a)$. 
We also assume that the transition density $p(s',\varepsilon'|s,\varepsilon,a)$ factors as $g(\varepsilon'|s')p(s'|s,a)$,
for a conditional density $g(\varepsilon'|s')$. If 
we assume $\varepsilon$ has a multivariate Gumbel or Type 1 extreme value distribution given by
\begin{equation}
    F(\varepsilon|s) = \prod_{a \in A(s)} \exp\{-\exp\{-(\varepsilon(a)-\mu)/\sigma\}\},
\label{eq:evd}
\end{equation}
where we set $\mu=-\sigma\gamma$ where $\gamma$ is Euler's constant, then $E\{\varepsilon(a)|s\}=0$ for $a \in A(s)$. 
Let $V(s,\varepsilon)$ be the optimal value of this MDP, the solution to
\begin{equation}
       V(s,\varepsilon)=\max_\pi  V_\pi(s,\varepsilon),  \quad \mbox{where} \;  V_\pi(s,\varepsilon) = E_\pi \left\{\sum_{t=0}^\infty \beta^t[r(s_t,a_t)+\varepsilon_t(a_t)]|s_0=s,\varepsilon_0=\varepsilon\right\}.
\label{eq:smoothVeps}
\end{equation}
Under these assumptions we can write the Bellman equation (\ref{eq:bellman}) for $V(s,\varepsilon)$ and show it 
has the representation
\begin{equation} V(s,\varepsilon)=\max_{a \in A(s)}[Q(s,a)+\epsilon(a)],
\label{eq:Vrep}
\end{equation}
where $Q=\Lambda_\sigma(Q)$ is the unique fixed point to the ``smoothed Bellman operator'' $\Lambda_\sigma$ given by
\begin{equation}
    \Lambda_\sigma(Q)(s,a) = r(s,a) + \beta \sum_{s'} \sigma \log\left(\sum_{a' \in A(s')} \exp\{Q(s',a')/\sigma\}\right)p(s'|s,a).
\label{eq:smoothQ}
\end{equation}
The derivation of (\ref{eq:Vrep}) and (\ref{eq:smoothQ}) use an important
property of the extreme value distribution for $\varepsilon$, namely that the expectation of $V(s,\varepsilon)$ with respect to $\varepsilon$
has an  analytical expression given by the so-called log-sum or ``smoothed max function'' $V(s)$ 
\begin{eqnarray}
    V(s) &\equiv &  E\left\{V(\tilde s,\tilde \varepsilon)\left|\tilde s=s\right\}\right. 
    =  \int_\varepsilon \max_{a \in A(s)} \left[Q(s,a)+\varepsilon(a)\right]F(d\varepsilon|s) \nonumber \\ 
    &=&  \sigma \log\left(\sum_{a\in A(s)} \exp\{Q(s,a)/\sigma\}\right).
\label{eq:smoothV}
\end{eqnarray}
Using equations (\ref{eq:smoothQ}) and (\ref{eq:smoothV}) we can derive a ``smoothed Bellman equation'' for $V$
\begin{equation}
    V(s)= \sigma \log\left(\sum_{a\in A(s)} \exp\left\{[r(s,a)+\beta \sum_{s'} V(s')p(s'|s,a)]/\sigma\right\}\right).
 \label{eq:smoothBellman}
\end{equation}
Equation (\ref{eq:smoothQ}) defines a ``smoothed $Q$'' as the fixed point of the smoothed contraction mapping 
$Q_\sigma=\Lambda_\sigma(Q_\sigma)$, and (\ref{eq:smoothBellman}) defines $V$ as a fixed point of the smoothed 
Bellman operator (also a contraction mapping)
$V_\sigma=\Gamma_\sigma(V_\sigma)$. These operators are equivalent  to the corresponding operators for MDP without the state variable
$\varepsilon$ in equations (\ref{eq:q1}) and (\ref{eq:q2}) and it is not hard to show that $Q_\sigma$ and $V_\sigma$ converge
to $Q$ and $V$ as $\sigma \downarrow 0$. 


The representation of $V(x,\varepsilon)$ in equation (\ref{eq:Vrep}) means that the decision rule $\pi(a|x,\varepsilon)$ 
is isomorphic to a static discrete choice model where $Q(s,a)+\varepsilon(a)$ is the utility/payoff from action $a$.
The CCP $\pi(a|s)$ is the conditional probability that the DM chooses action $a \in A(s)$
given the observed state $s$, taking the multinomial logit form \citep{mcfadden:1973}
\begin{eqnarray}
     \pi(a|s)  & =& \int I\left\{Q(s,a)+\varepsilon(a) \ge \max_{a' \in A(s)}[Q(s,a')+\varepsilon(a')]\right\}F(d\varepsilon|s)\nonumber \\
 &=&{ \exp\{Q(s,a)/\sigma\} \over \sum_{a' \in A(s)} \exp\{Q(s,a')/\sigma\}}.
\label{eq:ccp}
\end{eqnarray}
As $\sigma \downarrow 0$, $\pi(a|s)$ converges to a degenerate probability that equals 1 if 
$Q(s,a)$ is the largest $Q$ value  and 0 otherwise, i.e. to the optimal pure strategy $\pi$
 in equation (\ref{eq:bellman_policy}) for MDPs without the  $\varepsilon$ variable.
If follows that the choice-specific values $Q(s,a)$ in (\ref{eq:smoothQ}) and CCPs $\pi(a|s)$ in (\ref{eq:ccp})
are exactly the same as the ``soft Q'' values and optimal policy derived in the RL
literature under the assumption that agents care about entropy of the policy $\pi$, given in equation
(\ref{eq:vdefentropy}) of section 2.4.

Equation (\ref{eq:ccp}) can be considered as the analog of the {\it policy improvement\/} step for MDPs
in equation (\ref{eq:bellman_policy}).  The analog of the {\it policy valuation\/} step (\ref{eq:policy_valuation}) \citep{AguirregabiriaMira:2002} is 
the value function $V_\pi$  while is the solution to the linear system
\begin{equation}
   V_\pi(s) =\sum_{a \in A(s)} \pi(a|s)[r(s,a)+E\{\varepsilon(a)|s,a\}+\beta \sum_{s'} V_\pi(s')p(s'|s,a)],
\label{eq:smoothpolicyvaluation}
\end{equation}
where $E\{\varepsilon(a)|s,a\}$ is the conditional expectation of $\varepsilon(a)$ given that action $a$ is chosen
in state $s$  under decision rule $\pi$. The extreme value assumption implies that
$E\{\varepsilon(a)|s,a\}=-\sigma \log(\pi(a|s))$, and using equation (\ref{eq:ccp}), we see that right side of 
(\ref{eq:smoothpolicyvaluation}) equals the smoothed max function defining the smoothed Bellman operator in equation
(\ref{eq:smoothBellman}). It follows that we can use policy iteration to solve the smoothed Bellman equation $V_\sigma=\Gamma_\sigma(V_\sigma)$ 
cycling between policy valuation steps (\ref{eq:smoothpolicyvaluation})
and policy improvement steps (\ref{eq:ccp}) using the $Q_\pi$ implied by $V_\pi$ in equation (\ref{eq:smoothQ}).


Equation (\ref{eq:ccp}) implies that $\sigma\log(\pi(a|s))=Q(s,a)-V(s)\equiv \mathcal{A}(s,a)$, where
$\mathcal{A}(s,a)$ is called the {\it advantage function\/} in RL. Using $\mathcal{A}(s,a)$ we can rewrite the nonlinear fixed point
equation (\ref{eq:smoothQ}) for $Q$  as a linear system depending  on $r(s,a)$ and $\log(\pi(a|s))$,
\begin{equation}
    Q(s,a) = r(s,a) + \beta \sum_{s'} \sum_{a' \in A(s')} \left[ -\sigma\log(\pi(a'|s'))+Q(s',a')\right]\pi(a'|s')p(s'|s,a).
\label{eq:linq}
\end{equation}
Equation (\ref{eq:linq}) implies that we can decompose $Q$ as $Q(s,a)=Q_r(s,a)+Q_\varepsilon(s,a)$ where $Q_r(s,a)$ is the component 
of discounted value from observed rewards $r$ and $Q_\varepsilon(s,a)$ is the component from unobserved rewards $\varepsilon$.
Further substituting expression (\ref{eq:smoothpolicyvaluation}) in place of (\ref{eq:smoothBellman}) in the equation for
$Q$ in (\ref{eq:smoothQ}), we can show that $Q_r$ and $Q_\varepsilon$ are each the solution to the following linear systems
\begin{eqnarray}
     Q_r(s,a)&=&r(s,a)+\beta \sum_{s'} \sum_{a' \in A(s')} Q_r(s',a')\pi(a'|s')p(s'|s,a) \label{eq:Qr} \\
     Q_\varepsilon(s,a)&=&\beta \sum_{s'} \sum_{a' \in A(s')}\left[-\sigma\log(\pi(a'|s'))+ Q_\varepsilon(s',a')\right]\pi(a'|s')p(s'|s,a).
\label{eq:Qeps}
\end{eqnarray}
This decomposition reduces the computational burden of estimating DDC models.

\subsection{Estimation Methods for DDCs}
\label{section:ddc_methods}

The most common estimation method for DDCs is some form of  maximum likelihood.
Suppose we observe a random sample of {\it trajectories\/} i.e. data on the  consecutive states and actions taken by
 $N$  individuals $\{(s_{it},a_{it})|t=0,\ldots,T_i, \; i=1,\ldots,N\}$. Suppose  
that we smoothly parametrize $(r,p)$ using a vector unknown parameters $\theta$, which we denote by
$r_\theta$ and $p_\theta$.\footnote{\footnotesize
The discount factor  $\beta$ is often treated as known, though in 
some cases it is also estimated and thus part of $\theta$.}  
Using the implicit function theorem, we can show that $V_\theta$ and $Q_\theta$, the contraction fixed points given in equations
(\ref{eq:smoothV}) and (\ref{eq:smoothQ}), are continuously differentiable function of $\theta$.
Using these fixed points and the logit formula for $\pi(a|s)$ in equation (\ref{eq:ccp}) we can  write a full  likelihood
function for the data as a function of the unknown parameters $\theta$ given by
\begin{equation}
    L(\theta) = \prod_{i=1}^N \prod_{t=1}^{T_i}  \pi_\theta(a_{it}|s_{it})p_\theta(s_{it}|s_{it-1},a_{it-1}).
\label{eq:lf}
\end{equation}
The smoothness of $V_\theta$ and $Q_\theta$ in $\theta$ implies that the CCP $\pi_\theta$ and the
likelihood $L(\theta)$ are both smooth functions of $\theta$ so the standard asymptotic
efficiency properties of maximum likelihood apply. 

The Nested Fixed Point algorithm (NFXP) \citep{rust:1988} computes the MLE of $\theta$ via a standard outer hill-climbing loop that maximizes $L(\theta)$, where each evaluation calculates
the contraction fixed point $Q_\theta=\Lambda_\theta(Q_\theta)$ that defines the choice-specific values $Q_\theta$ where $\Lambda_\theta$
is the operator defined on the right hand side of equation (\ref{eq:smoothQ}).
Notice that $L(\theta)$ depends on $\theta$ via $Q_\theta$. An alternative strategy \citep{sujudd:2012}
computes the MLE via the MPEC algorithm, maximizing $L(\theta,Q)$ as a function of $(\theta,Q)$
subject to the fixed point constraint $Q_\theta=\Lambda_\theta(Q_\theta)$.\footnote{\footnotesize \citet{sujudd:2012} claim that the
MPEC alogithm is substantially faster than NFXP, but this was due to using SA \citep{sujuddcomment:2016}
to compute the fixed point $Q_\theta=\Lambda_\theta(Q_\theta)$, which is slow when $\beta$ is close to 1. When Newton's method is used
to compute this fixed point, the cpu times for MPEC and NFXP are roughly comparable with NFXP being faster on test problems with
large sample sizes.}
NFXP's computational burden stems from computing the inner fixed point $Q_\theta=\Lambda_\theta(Q_\theta)$
for each trial value of $\theta$ in the outer hill-climbing algorithm. 
Newton's method (e.g. policy iteration) is typically used to solve for $Q_\theta$ (or alternatively for $V_\theta$ in equation
(\ref{eq:smoothpolicyvaluation})). Even though  Newton's method converges rapidly, it requires $O(|S|^3)$ operations per Newton/policy
iteration step, making NFXP burdensome for large state spaces $S$. Therefore alternative estimation methods have been proposed that 
reduce the total number of policy iteration steps required to estimate $\theta$.


The ``CCP estimator'' \citep{hotzmiller:1993}
for $r$ uses a non-parametric estimate of $\pi(a|s)$ to avoid repeated
solution of $Q$ over the search for $\theta$ that NFXP requires. Consider the case where $r_\theta$ is linear in parameters,
i.e. we can write $r_\theta(s,a)=r(s,a)*\theta \equiv \sum_{k=1}^K r_k(s,a)\theta_k$ for known functions
 $\vec{r}=(r_1,\ldots,r_K)$ called {\it features\/} in the IRL
literature. Then (\ref{eq:Qr}) implies that $Q_r(s,a)=R(s,a)*\theta$ where
$R$ is the solution to
\begin{equation}
   R(s,a)=\vec{r}(s,a) + \beta \sum_{s'} \sum_{a' \in A(s')} R(s',a')\pi(a'|s')p(s'|s,a) = \vec{r}(s,a)+\beta ER(s,a),
\label{eq:hequation}
\end{equation}
where $ER(s,a)$ is a shorthand for the conditional expectation of $R(s,a)$ the middle of expression in (\ref{eq:hequation}).
It follows that we only need to solve (\ref{eq:hequation}) once using a non-parametric estimate of $\hat\pi$ along with
$Q_\varepsilon$ in (\ref{eq:Qeps}) to serve as a ``correction term''  to obtain the following form for the CCP
\begin{equation}  
           \pi_\theta(a|s) = {\exp\{\hat R(s,a)\cdot\theta+\hat Q_\varepsilon(s,a)\} \over \sum_{a' \in A(s)}
\exp\{\hat R(s,a')\cdot\theta+\hat Q_\varepsilon(s,a')\}},
\label{eq:ccpestimator}
\end{equation}
where $\hat R(s,a)$ is the solution to (\ref{eq:hequation}) and $\hat Q_\varepsilon$ is the solution
to (\ref{eq:Qeps}) using non-parametric first stage
estimators $\hat\pi(a|s)$ and $\hat p(s'|s,a)$. The CCP estimator $\hat\theta$ maximizes the partial likelihood $L_\pi(\theta)$ given by 
\begin{equation}
    L_\pi(\theta) = \prod_{i=1}^N \prod_{t=1}^{T_i}  \pi_\theta(a_{it}|s_{it}),
\label{eq:ccplf}
\end{equation}
where $\pi$ is given in (\ref{eq:ccpestimator}).\footnote{\footnotesize See \citet{lrspe:2022} 
for  a modified version of the CCP estimator that is ``locally robust'' i.e. it corrects for sampling
error in the first stage estimates of the functions $\hat H(s,a)$ and $Q_\varepsilon(s,a)$.}
Thus, when $r$ is a  linear combination of known features, we only  need to compute $\hat R$ and $\hat Q_\varepsilon$ once
at the start of the estimation, requiring  $K+1$ solutions of systems with $|S|$ equations and unknowns,
where $K$ is the number of features. This can be further reduced to
a single solution if we assume that $r(s,a_s)$ is known for some
``normalizing alternative'' $a_s \in A(s)$ (see Section~\ref{section:identification}).

A cost of the CCP estimator is that it is less efficient than using NFXP to estimate the same partial likelihood criterion (\ref{eq:ccplf})
due to the noise in the first stage non-paramteric estimates $\hat\pi$ that are used to construct
the $\hat Q_\varepsilon$ function in (\ref{eq:ccpestimator}).  \citet{AguirregabiriaMira:2002} (AM) addressed this problem
by proposing an iterative estimator called {\it nested pseudo likelihood\/} (NPL)  that uses an initial non-parametric estimate $\hat\pi$ 
to compute the value function $V_{\hat\pi}$ implied by this initial estimate by solving for it in the policy valuation step
in equation (\ref{eq:smoothpolicyvaluation}), which we denote as $\hat V_{\theta,\hat\pi}$ since it depends
both on $\hat\pi$ and the parameters $\theta$ entering
the reward function. AM then  estimate  $\hat\theta$ using a pseudo-likelihood similar to (\ref{eq:ccplf})
except that the CCP $\pi_\theta(a|s)$ is evaluated using (\ref{eq:ccp}) using estimated $Q$ functions given by
\begin{equation}
                  \hat Q_\theta(s,a)=r_\theta(s,a)+\beta \sum_{s'}\hat V_{\theta,\hat\pi}(s'|s,a)p(s'|s,a),
\label{eq:qhat}
\end{equation}
Thus, the likelihood $L_\pi(\theta)$ as depends implicitly on the first stage non-parametric CCP estimate $\hat\pi$, which can be regarded
as a ``nuisance parameter''. In principle, estimation noise in $\hat\pi$ will contaminate and affect the asymptotic 
covariance matrix for the parameters of interest $\theta$. However 
AM established a key {\it zero Jacobian property\/} that at the fixed point $\partial V(s)/\partial \pi=0$ 
which  implies that  the first stage nuisance parameter $\hat\pi$
is ``Neyman orthogonal'' with respect to the parameters of interest $\theta$ in the sense of \citet{lrspe:2022}. It follows 
that the asymptotic covariance matrix
for the NPL estimator of $\theta$ is unaffected by noise in the first stage $\hat\pi$ and is the same 
as $\pi$  was known. Further, 
AM showed that NPL can be iterated: using $\hat\theta$, they can compute an updated estimate of the CCPs 
$\hat\pi=\pi_{\hat\theta})$ and repeat the policy valuation step to get a new set of $Q$ functions (\ref{eq:qhat}), and use 
these in the likelihood (\ref{eq:ccplf}) to get an updated estimate $\hat\theta$. They refer to NPL as ``swapping''
the policy iteration and maximization steps of the NFXP estimator (which does maximization in the ``outside'' loop
and policy iteration in the ``inside'' or nested fixed point loop). Swapping the order of maximization
and policy iteration reduces the total number of policy iteration steps required to 
estimate $\theta$. AM showed that if this iterative version of  NPL converges, it produces the
same  estimate $\hat\theta$ that  NFXP produces from maximization of the partial likelihood  (\ref{eq:ccplf}).
Thus the NPL estimator is as efficient as NFXP but at lower computational cost, and it is more efficient than the CCP
estimator but at a higher computational cost.\footnote{\footnotesize
See also {\it efficient pseudo likelihood\/} (EPL) \citep{dearingblevins:2025}, a similar iterative
procedure that is also an efficient sequential estimator that 
can estimate parameters of dynamic games.}  

Bayesian approaches have been proposed for inference in DDC models using Markov Chain Monte Carlo methods that simultaneously
simulate the likelihood as well as stochastic simulations to compute estimates of value functions that are similar to
``Monte Carlo tree search'' used by RL algorithms such as the Q-learning algorithm discussed in section 2.4, 
see \citet{ijc:2009}. There are also non-Bayesian simulation estimators such as
\citet{bbl:2007} that estimate $\theta$ as the value that minimizes the degree of suboptimality
in observed actions relative to other actions not taken  using simulation estimates of $Q$ function evaluated at the actions
the DM actually took and other actions not chosen.

\subsection{Function Approximation Methods for Large State Spaces}

When the state space $S$ is large (i.e. there are continuous state variables that must be discretized)
the approaches discussed above that use PI and repeated solution of 
large linear systems with $|S|$ equations and unknowns become computationally infeasible. 
Following section~\ref{section:curse} we can approximate 
the fixed point $Q=\Lambda_\sigma(Q)$ of the smoothed Bellman operator in equation (\ref{eq:smoothQ}) by $Q_\phi(s)$
with a flexible parametric family such as a neural network with weights $\phi$, and solve 
the fixed point problems (\ref{eq:smoothQ}) and (\ref{eq:smoothV}) more tractably as nonlinear least squares problems.  
Let $\rho(\theta,\phi)$ be the sum of squared residuals of the smoothed Bellman operator, $Q_\phi-\Lambda_\sigma(Q_\phi)$,
over a grid of $N$ points in the state space, equation (\ref{eq:dqnideal}). Note that the residuals also depend on $\theta$ via $\{\beta,r,p\}$
in the smoothed Bellman operator $\Lambda_\sigma$.
A nested least squares estimator similar to NFXP can estimate $\theta$ 
by maximizing the likelihood $L(\theta)$ (\ref{eq:lf}) while mimimizing $\rho(\theta,\phi)$ over $\phi$ for
each value of $\theta$. A penalized likelihood estimator called SEES \citep{LuoSang:2024}, for {\it sieve-based efficient estimator for
structural models\/}, chooses $\hat\theta$ and $\hat\phi$ to maximize 
\begin{equation}
(\hat\theta,\hat\phi)=\argmax_{\theta,\phi} \log(L(\theta))-\lambda \rho(\theta,\phi),
\label{eq:penalizedllf}
\end{equation}
where $\lambda \ge 0$ is a penalization parameter. Consistency of SEES requires the use of a {\it sieve,\/} i.e.
the dimension of $\phi$ must grow to infinity at certain rate with the sample
size so that the approximation error in approximating $Q$ by $Q_\phi$ converges to zero.
A related estimator called {\it neural network efficient estimator\/} (NNES) \citep{Nguyen:2025} maximizes a penalized criterion
similar to the SEES estimator (\ref{eq:penalizedllf}) but in a sequential fashion similar to the NPL
estimator of AM.

A challenge facing sieve-based estimators is that $\phi$ are ``nuisance parameters'' 
(their only purpose is to approximate $Q$) whereas $\theta$ are the ``parameters of interest'' 
that determine  rewards and beliefs, $\{\beta,r,p\}$. Since $\hat\phi$  minimizes
the sum of squared residuals $\rho(\phi,\theta)$  for each $\theta$, it follows the $\hat\phi$ will be an implicit
function of $\theta$. The asymptotic covariance matrix for $\hat\theta$  requires $\nabla_\theta Q$, the gradient
of $Q$ with respect to $\theta$. However directly differentiating
$Q_{\hat\phi(\theta)}$ with respect to $\theta$ using the chain rule
 $\nabla_\theta Q=\nabla_\phi Q_{\hat\phi(\theta)}\nabla_\theta\hat\phi(\theta)$ yields inaccurate results due to the
high dimensionality of $\phi$ and the fact that the minimization over  $\phi$ can result in local optima and fragile
solutions that can cause small changes in $\theta$ to produce large changes in $\hat\phi(\theta)$. 
\citet{Nguyen:2025} addresses this problem by 
differentiating both sides of the fixed point equation for $Q$ (\ref{eq:linq}) with respect to $\theta$ 
to obtain a separate fixed point equation that uniquely determines  $\nabla_\theta Q$. 
He approximates this solution by parametrizing $\nabla_\theta Q$ by a separate neural net with 
weights $\psi$, and shows that the solution $\nabla_\theta Q_{\hat\psi}$ accurately approximates
to $\nabla_\theta Q$. However doing this involves the computational burden of having to fit
two neural nets each time $\theta$ is updated: one for $Q_{\hat\phi(\theta)}$ and the other for $\nabla_\theta Q_{\hat\psi(\theta)}$.





\subsection{The Identification Problem}
\label{section:identification}

Section~\ref{sec:mdp} showed how $\pi$ is calculated from the structure $\{\beta,r,p\}$, which
we summarize by the mapping $\pi=f(\beta,r,p)$. The Identification Problem asks if this mapping is invertible:
given $\pi$ is there a unique underlying structure $\{\beta,r,p\}$ that implies it? Unfortunately, without
further restrictions the answer is negative, as shown by \citet{hotzmiller:1993}, 
\citet{rust:1994} and \citet{magnacthesmar:2002}. 
To understand why,  consider the first step of the inversion process, uncovering  $Q$ from $\pi$.
Using the logit form of $\pi$ in (\ref{eq:ccp}) we have
\begin{equation}
               \log\left(\pi(a|s)/\pi(a'|s)\right) = Q(s,a)-Q(s,a'), \quad s \in S, \; a,a' \in A.
\label{eq:ccpinverse}
\end{equation}
We already see that information on $\pi$ is insufficient to recover all the $Q$ values: only the {\it differences\/} in
$Q$ are identified but  not all $|S||A|$ possible values of $Q(s,a)$. Indeed, the structure $\{\beta,r,p\}$ contains $1+|S|(|A|-1)+|S|^2|A|$  
parameters which is more than then $|S|(|A|-1)$ parameters of $\pi$, so we will generally have ``more equations than unknowns'' 
and the underlying structure will only be partially identified.

Identification requires additional restrictions on $\{\beta,r,p\}$. 
Empirical studies commonly impose the assumption that $r$ and $p$ have known parametric functional forms that depend on a relatively small number of parameters $\theta$. In some situations the reward function can be treated as
known, and this facilitates the identification of subjective beliefs $p(s'|s,a)$.\footnote{\footnotesize 
For example \citet{tennis:2025} are able to identify the subjective beliefs of professional tennis servers about how their serve
strategies affects their probability of winning a tennis game. These beliefs take the form of a transition $p(s'|s,a)$ where $s$ denotes
the point state of the game and $a$ denotes serve direction. Identification of beliefs is possible  because it is reasonable to assume
$r$ is known: i.e. the server earns a reward of 1 if they win the game and 0 otherwise, and due to reasonable parametric restrictions
on $p$.}
In many applications researchers assume agents have rational expectations, 
i.e. their subjective beliefs about state transitions $p(s'|s,a)$ correspond to the actual probability governing
these transitions. With sufficient data, $p$ can be estimated non-parametrically and so for the analysis of identification
it can be treated as known.\footnote{\footnotesize With sufficient restrictions on $r$ the rational expectations
 assumption is testable. For example, \citet{tennis:2025} strongly reject the hypothesis the many professional tennis servers have rational beliefs about their
own ability and those of their opponent, as embodied by $p(s'|s,a)$.}
Many studies typically assume that the discount factor $\beta$ is also known, however under the assumption of 
rational expectations and restrictions on rewards, $\beta$ can be identified, see \citet{abbringdaljord:2020}.

If $\{\beta,p\}$ are known, identification of $r$ depends on the assumption that either it has a known parametric
functional form $r_\theta$ where the dimension of $\theta$ 
is less than $|S|(|A|-1)$, or that for each $s \in S$ there is
a normalizing or ``anchor action'' $a_s$ such that  $r(s,a_s)$ is known. The latter assumption amounts to $|S|$ restrictions
that imply there are only $|S|(|A|-1)$ unknown rewards to be recovered, the same number of free values of $\pi$. 
\citet{hotzmiller:1993} and \citet{kang:2025} show that in this ``exactly identified'' case,
we can calculate $Q(s,a_s)$ from knowledge of $r(s,a_s)$ and $\pi(a_s|s)$ using equation (\ref{eq:linq}) 
and then (\ref{eq:ccpinverse}) allows us to recover all remaining $Q$ values. 
Using the $Q$ values we can back out the remaining unknown rewards $r$ from the smooth Bellman equation (\ref{eq:smoothQ}).

The assumption that we know $r(s,a_s)$ for some anchor action $a_s$ in each state $s \in S$ is far from innocuous. If wrong, estimated rewards will not equal true rewards even though the model perfectly fits $\pi$, and counterfactual predictions may differ from actual outcomes. In the absence of prior restrictions on rewards, we can only identify ``observationally equivalent'' rewards of the form $r(s,a)+h(s)-\beta \sum_{s'} h(s')p(s'|s,a)$ for any function $h: S \to R$. This partially identified set defines an inherent limit on what can be learned about rewards without accurate prior knowledge of at least some aspects of the reward function.\footnote{\footnotesize This is what the IRL literature, following \citet{nhr:1999}, calls potential-based reward shaping, where $h(s)$ is a `potential function' $\Phi(s)$. Any reward $r'(s,a) = r(s,a) + \beta\Phi(s') - \Phi(s)$ yields the same optimal policy: identification `up to policy invariance.' For imitation learning, policy invariance suffices. But for structural econometrics, policy invariance is not enough.}

To resolve the identification ambiguity, the DDC and IRL literatures have developed distinct sets of identifying assumptions, each reflecting the typical problems and goals of the respective fields. \citet{souza:2021} provide examples of certain counterfactuals that are  identified even when $r$ is
only partially identified. They conclude that their ``results call for caution while leaving room for optimism: 
although counterfactual behavior and welfare can be sensitive to identifying restrictions imposed on the model, 
there exists important classes of counterfactuals that are robust to such restrictions.'' (p. 385).




